% GENERATED FILE. Edit canonical lessons, then run npm run generate:books.
% canonical-source-hash: fnv1a64:8b85a314ef52ee55
% canonical-lessons: TA-C538-pinnu, TA-C538-mutai, TA-C538-kor, TA-C538-tari, TA-C538-pattu

\chapter{Threads and Beads}
\label{ch:ta-threads-and-beads}
\hlchaptermodality{538}

\begin{chapteropening}
\textbf{By the end of this chapter:} \emph{I can talk about plaiting, weaving and stringing.}

Complete the last lesson of chapter 538: I can talk about plaiting, weaving and stringing.
\end{chapteropening}

\section[\texorpdfstring{\ta{பின்னு} (piṉṉu)}{பின்னு (piṉṉu)}]{\ta{பின்னு} (piṉṉu) — to knit; to plait}

\label{lesson:TA-C538-pinnu}



\begin{quote}
Before the new one: say the Tamil for to polish, then the Tamil for to bake.
\end{quote}

\subsection*{You'll want to know: \ta{பின்னு}}
\textbf{\ta{பின்னு}} — \emph{piṉṉu} — \textquotedblleft{}to knit; to plait\textquotedblright{}.

\begin{quote}
\ta{அம்மா} \ta{என்} \ta{தலைமுடியைப்} \ta{பின்னினாள்}.
\end{quote}

\begin{quote}
\emph{ammā eṉ talaimuṭiyaip piṉṉiṉāḷ.}
\end{quote}

\begin{quote}
Mother plaited my hair.
\end{quote}

Past (I …): \ta{பின்னினேன்}, \emph{piṉṉiṉēṉ}.

Of hair, wool or a net. The noun is \ta{பின்னல்}, a plait.

\begin{cousinweb}
\textbf{\ta{பின்னு}} (\emph{piṉṉu}) is native Dravidian: DEDR 4207 gives \ta{பின்னு} (\emph{piṉṉu}), \textquotedblleft{}to plait, to braid, to lace, to knit\textquotedblright{}.
\end{cousinweb}

\begin{grammarlens}[title={What you've built}]
One more word for describing what you lived and felt.
\end{grammarlens}

\subsection*{Your turn}
\begin{itemize}
\raggedright
  \item \emph{Say it:} \emph{piṉṉu}
  \item \emph{Say it:} \emph{piṉṉu}, once more
  \item \emph{Say it:} \emph{ammā eṉ talaimuṭiyaip piṉṉiṉāḷ.} — Mother plaited my hair.
  \item \emph{Recall:} say the Tamil for to polish, then the Tamil for to bake, then say \emph{piṉṉu} again
\end{itemize}

\begin{tcolorbox}[breakable,colback=teal!4,colframe=teal!35!black,title={Before you move on}]
What does \ta{பின்னு} mean? (\textquotedblleft{}To knit; to plait\textquotedblright{}.) Say it once more.
\end{tcolorbox}
\section[\texorpdfstring{\ta{முடை} (muṭai)}{முடை (muṭai)}]{\ta{முடை} (muṭai) — to weave (a basket or mat)}

\label{lesson:TA-C538-mutai}



\begin{quote}
Before the new one: say the Tamil for to bake, then the Tamil for to knit.
\end{quote}

\subsection*{You'll want to know: \ta{முடை}}
\textbf{\ta{முடை}} — \emph{muṭai} — \textquotedblleft{}to weave (a basket or mat)\textquotedblright{}.

\begin{quote}
\ta{பாட்டி} \ta{பனை} \ta{ஓலையால்} \ta{ஒரு} \ta{கூடை} \ta{முடைந்தாள்}.
\end{quote}

\begin{quote}
\emph{pāṭṭi paṉai ōlaiyāl oru kūṭai muṭaintāḷ.}
\end{quote}

\begin{quote}
Grandmother wove a basket from palm leaves.
\end{quote}

Past (I …): \ta{முடைந்தேன்}, \emph{muṭaintēṉ}.

Of baskets and mats; cloth is woven with \ta{நெய்} (\emph{ney}).

\begin{cousinweb}
\textbf{\ta{முடை}} (\emph{muṭai}) is native Dravidian: DEDR 4927 gives \ta{முடை} (\emph{muṭai}), \textquotedblleft{}to braid, to plait, to wattle\textquotedblright{}, and the noun for a basket of palm leaf.
\end{cousinweb}

\begin{grammarlens}[title={What you've built}]
One more word for describing what you lived and felt.
\end{grammarlens}

\subsection*{Your turn}
\begin{itemize}
\raggedright
  \item \emph{Say it:} \emph{muṭai}
  \item \emph{Say it:} \emph{muṭai}, once more
  \item \emph{Say it:} \emph{pāṭṭi paṉai ōlaiyāl oru kūṭai muṭaintāḷ.} — Grandmother wove a basket from palm leaves.
  \item \emph{Recall:} say the Tamil for to bake, then the Tamil for to knit, then say \emph{muṭai} again
\end{itemize}

\begin{tcolorbox}[breakable,colback=teal!4,colframe=teal!35!black,title={Before you move on}]
What does \ta{முடை} mean? (\textquotedblleft{}To weave (a basket or mat)\textquotedblright{}.) Say it once more.
\end{tcolorbox}
\section[\texorpdfstring{\ta{கோர்} (kōr)}{கோர் (kōr)}]{\ta{கோர்} (kōr) — to string (beads, flowers)}

\label{lesson:TA-C538-kor}



\begin{quote}
Before the new one: say the Tamil for to knit, then the Tamil for to weave.
\end{quote}

\subsection*{You'll want to know: \ta{கோர்}}
\textbf{\ta{கோர்}} — \emph{kōr} — \textquotedblleft{}to string (beads, flowers)\textquotedblright{}.

\begin{quote}
\ta{அவள்} \ta{மல்லிகைப்} \ta{பூக்களை} \ta{நூலில்} \ta{கோர்த்தாள்}.
\end{quote}

\begin{quote}
\emph{avaḷ mallikaip pūkkaḷai nūlil kōrttāḷ.}
\end{quote}

\begin{quote}
She strung jasmine flowers on a thread.
\end{quote}

Past (I …): \ta{கோர்த்தேன்}, \emph{kōrttēṉ}.

What is strung takes -\ta{ஐ}, the thread -\ta{இல்}: \ta{பூக்களை} \ta{நூலில்} \ta{கோர்}.

\begin{cousinweb}
\textbf{\ta{கோர்}} (\emph{kōr}) is native Dravidian: DEDR 2176 gives \ta{கோ} and \ta{கோர்} (\emph{kōr}), \textquotedblleft{}to string beads or flowers, to thread a needle\textquotedblright{}.
\end{cousinweb}

\begin{grammarlens}[title={What you've built}]
One more word for describing what you lived and felt.
\end{grammarlens}

\subsection*{Your turn}
\begin{itemize}
\raggedright
  \item \emph{Say it:} \emph{kōr}
  \item \emph{Say it:} \emph{kōr}, once more
  \item \emph{Say it:} \emph{avaḷ mallikaip pūkkaḷai nūlil kōrttāḷ.} — She strung jasmine flowers on a thread.
  \item \emph{Recall:} say the Tamil for to knit, then the Tamil for to weave, then say \emph{kōr} again
\end{itemize}

\begin{tcolorbox}[breakable,colback=teal!4,colframe=teal!35!black,title={Before you move on}]
What does \ta{கோர்} mean? (\textquotedblleft{}To string (beads, flowers)\textquotedblright{}.) Say it once more.
\end{tcolorbox}
\section[\texorpdfstring{\ta{தறி} (taṟi)}{தறி (taṟi)}]{\ta{தறி} (taṟi) — a loom}

\label{lesson:TA-C538-tari}



\begin{quote}
Before the new one: say the Tamil for to weave, then the Tamil for to string.
\end{quote}

\subsection*{You'll want to know: \ta{தறி}}
\textbf{\ta{தறி}} — \emph{taṟi} — \textquotedblleft{}a loom\textquotedblright{}.

\begin{quote}
\ta{நெசவாளர்} \ta{தறியில்} \ta{சேலை} \ta{நெய்கிறார்}.
\end{quote}

\begin{quote}
\emph{necavāḷar taṟiyil cēlai neykiṟār.}
\end{quote}

\begin{quote}
The weaver weaves a sari on the loom.
\end{quote}

A hand loom is a \ta{கைத்தறி}.

\begin{cousinweb}
\textbf{\ta{தறி}} (\emph{taṟi}) is native Dravidian: DEDR 3140 gives \ta{தறி} (\emph{taṟi}), \textquotedblleft{}to cut off\textquotedblright{}, and the noun for a wooden post and a weaver's loom.
\end{cousinweb}

\begin{grammarlens}[title={What you've built}]
One more word for describing what you lived and felt.
\end{grammarlens}

\subsection*{Your turn}
\begin{itemize}
\raggedright
  \item \emph{Say it:} \emph{taṟi}
  \item \emph{Say it:} \emph{taṟi}, once more
  \item \emph{Say it:} \emph{necavāḷar taṟiyil cēlai neykiṟār.} — The weaver weaves a sari on the loom.
  \item \emph{Recall:} say the Tamil for to weave, then the Tamil for to string, then say \emph{taṟi} again
\end{itemize}

\begin{tcolorbox}[breakable,colback=teal!4,colframe=teal!35!black,title={Before you move on}]
What does \ta{தறி} mean? (\textquotedblleft{}A loom\textquotedblright{}.) Say it once more.
\end{tcolorbox}
\section[\texorpdfstring{\ta{பட்டு} (paṭṭu)}{பட்டு (paṭṭu)}]{\ta{பட்டு} (paṭṭu) — silk}

\label{lesson:TA-C538-pattu}



\begin{quote}
Before the new one: say the Tamil for to string, then the Tamil for a loom.
\end{quote}

\subsection*{You'll want to know: \ta{பட்டு}}
\textbf{\ta{பட்டு}} — \emph{paṭṭu} — \textquotedblleft{}silk\textquotedblright{}.

\begin{quote}
\ta{திருமணத்துக்கு} \ta{அவள்} \ta{பட்டுச்} \ta{சேலை} \ta{அணிந்தாள்}.
\end{quote}

\begin{quote}
\emph{tirumaṇattukku avaḷ paṭṭuc cēlai aṇintāḷ.}
\end{quote}

\begin{quote}
She wore a silk sari for the wedding.
\end{quote}

Before a noun it doubles the next stop: \ta{பட்டுச்} \ta{சேலை}, a silk sari.

\begin{cousinweb}
\textbf{\ta{பட்டு}} (\emph{paṭṭu}) is borrowed from Sanskrit \emph{paṭṭa}, \textquotedblleft{}cloth, woven silk\textquotedblright{} (CDIAL 7700), with the Tamil final -\ta{உ}.
\end{cousinweb}

\begin{grammarlens}[title={What you've built}]
One more word for describing what you lived and felt.
\end{grammarlens}

\subsection*{Your turn}
\begin{itemize}
\raggedright
  \item \emph{Say it:} \emph{paṭṭu}
  \item \emph{Say it:} \emph{paṭṭu}, once more
  \item \emph{Say it:} \emph{tirumaṇattukku avaḷ paṭṭuc cēlai aṇintāḷ.} — She wore a silk sari for the wedding.
  \item \emph{Recall:} say the Tamil for to string, then the Tamil for a loom, then say \emph{paṭṭu} again
\end{itemize}

\begin{tcolorbox}[breakable,colback=teal!4,colframe=teal!35!black,title={Before you move on}]
What does \ta{பட்டு} mean? (\textquotedblleft{}Silk\textquotedblright{}.) Say it once more.
\end{tcolorbox}
